\CUSTOMstruct{ВВЕДЕНИЕ}

Производственная практика является обязательным этапом образовательного процесса, который направлен на проверку компетенцией обучающегося.
Данная проверка подразумевает по собой проверку компетенций, полученных в процессе освоения образовательной программы направления -- системное прикладное и программное обеспечение, путём работы над настоящими задачами в настоящем проекте, 
где трудовой распорядок максимально приближен к официальной занятости посредством трудоустройства.

Данный отчёт отражает работу в рамках практики. 

Цель практики заключалась в демонстрации учащимися приобретённых навыков работы по специальности.

Чтобы выполнить поставленную цель были проделаны следующие этапы, приведённые в таблице 0.1.

\begin{table}[h]
  \centering
  \caption{Этапы и задания практики}
  \label{tab:practice-stages}
  \renewcommand{\arraystretch}{1.2}
  \begin{tabularx}{\textwidth}{|>{\raggedright\arraybackslash}p{0.28\textwidth}|>{\raggedright\arraybackslash}X|}
    \hline
    \textbf{Название этапа} & \textbf{Задание} \\
    \hline
    Инструктаж обучающегося & Инструктаж обучающегося по ознакомлению с требованиями охраны труда, техники безопасности, пожарной безопасности, а также правилами внутреннего трудового распорядка \\
    \hline
    Утилита для сбора логов по camunda events & - Сервис реализует API для получения истории событий Camunda за произвольный временной диапазон с параметрами startTime и endTime, передаваемыми через URI. - Возвращаемые данные включают только активности типов: service task, user task, exclusive gateway, send task, start event, script task, inclusive gateway, intermediate message throw event. - Формат выгрузки данных — .xlsx, со столбцами Activity ID, Activity Name, Start Time, End Time, Duration (ms), Activity Type. - Сбор истории выполняется без заметного влияния на производительность основной системы (разница в средней загрузке CPU не более 5\%). - Ошибки при запросе возвращают корректный HTTP-ответ с кодами: 400 (некорректные параметры), 500 (внутренняя ошибка). В отчёт войдёт: описание шагов для аналитики решения; описание шагов решения; само решение в коде (контроллер, сервис, возможно затронутое api и его дописывание); файлы, которые будут формироваться при запросе; обработка ошибок; замер метрик цпу при нагрузке на камунду для сверки средней нагрузки (проверка интрузивности); вывод. \\
    \hline
    Дописывание получения движений токена по workflow & - Получение истории и статуса токенов переведено на асинхронную схему: Camunda самостоятельно публикует уведомления о событиях в брокер сообщений. - На стороне бизнес-логики реализован модуль подписки на эти уведомления без периодического опроса (polling). - Подключение брокера выполнено в виде плагина Camunda Engine, активируемого при старте standalone-сборки. - Система корректно восстанавливает состояние подписки после перезапуска. \\
    \hline
  \end{tabularx}
\end{table}

\newpage
\begin{table}[h]
  \centering
  \caption*{Продолжение таблицы~\ref{tab:practice-stages}}
  \renewcommand{\arraystretch}{1.2}
  \begin{tabularx}{\textwidth}{|>{\raggedright\arraybackslash}p{0.28\textwidth}|>{\raggedright\arraybackslash}X|}
    \hline
    \textbf{Название этапа} & \textbf{Задание} \\
    \hline
    & В отчёт войдёт: описание шагов для аналитики решения; описание шагов решения; само решение в коде (хендлер, сервис/ы, затронутое api, его дописывание и код плагина); результаты сообщений из брокера; проверка корректного восстановления и перечитывания пропущенных событий из брокера; вывод. \\
    \hline
    Настройка движка Camunda под нагрузку & - Настроен встроенный брокер Camunda для эффективной работы при высокой нагрузке с external tasks. - Job Executors оптимизированы под асинхронное выполнение задач с минимальным временем ожидания. - Запись истории в слой персистенции отключена; при этом логика работы брокера полностью сохраняется. В отчёт войдёт: описание шагов для аналитики решения; описание шагов решения; само решение (конфигурационные файлы); проверка быстрой асинхронной работы job executors; проверка отсутствия истории в слое персистенции при движении по схеме; проверка, что это решение не затронуло брокер; вывод. \\
    \hline
    Написание кеша для метаданных & - Реализован read model кеш, наполняемый при старте сервиса из слоя персистенции. - Кеш подключается как отдельная библиотека с минимальным набором зависимостей. - Инвалидация кеша осуществляется по событиям через event listener без перезапуска сервиса. - При обращении к данным из кеша количество обращений к базе данных происходит только на заполнение кеша. - Кеш корректно обновляется при получении события изменения метаданных. В отчёт войдёт: описание шагов для аналитики решения; описание шагов решения; само решение кодом (сервисы библиотеки, зависимости, интеграция библиотеки); проверка заполенения кеша; проверка чтения из кеша; проверка инвалидации кеша; проверка снижения обращений в базу (cache hit rate); проверка обновлений кеша; вывод. \\
    \hline
  \end{tabularx}
\end{table}

\newpage
\begin{table}[h]
  \centering
  \caption*{Продолжение таблицы~\ref{tab:practice-stages}}
  \renewcommand{\arraystretch}{1.2}
  \begin{tabularx}{\textwidth}{|>{\raggedright\arraybackslash}p{0.28\textwidth}|>{\raggedright\arraybackslash}X|}
    \hline
    \textbf{Название этапа} & \textbf{Задание} \\
    \hline
    Оформление отчётных документов в соответствии с требованиями & 1. Должно быть подробное описание выполнения задач по этапам. Результаты задания необходимо разместить в приложения. 2. Оформление отчёта должно быть выполнено в соответствии с методическим пособием (https://books.ifmo.ru/file/pdf/2622.pdf) 3. Структура документа: титульный лист, введение, основная часть, заключение, приложения. 4. В основной части подробно описывается выполнение задач 2-5 этапов, в приложении помещаются результаты данных этапов. 5. Отчёт необходимо подгрузить в модуле практика как "письменный отчёт". \\
    \hline
    Получение отзыва руководителя практики & Получить отзыв с оценкой у руководителя практики в модуле Практики. \\
    \hline
    Техническое задание по ВКР & Написать техническое задание по ВКР. Техническое задание должно содержать: наименование, назначение, основание для разработки, функции, структура, пользовательский интерфейс, требования к надёжности, требования к безопасности, условия эксплуатации и проч.\ важные требования, документация, стадии и этапы разработки, порядок контроля и приёмка. В случае отсутствия какого-либо пункта необходимо обосновать его отсутствие. Сформировать в виде документа в формате pdf и прислать на проверку на электронную почту преподавателя. \\
    \hline
    Оформление отчётных документов в соответствии с требованиями & 1.~Должно быть подробное описание выполнения задач по этапам. Результаты задания необходимо разместить в приложения. 2.~Оформление отчёта должно быть выполнено в соответствии с методическим пособием (\url{https://books.ifmo.ru/file/pdf/2622.pdf}). 3.~Структура документа: титульный лист, введение, основная часть, заключение, приложения. 4.~В основной части подробно описывается выполнение задач 2--4 этапов, в приложении помещаются результаты данных\\
    \hline
  \end{tabularx}
\end{table}

\newpage
\begin{table}[h]
  \centering
  \caption*{Продолжение таблицы~\ref{tab:practice-stages}}
  \renewcommand{\arraystretch}{1.2}
  \begin{tabularx}{\textwidth}{|>{\raggedright\arraybackslash}p{0.28\textwidth}|>{\raggedright\arraybackslash}X|}
    \hline
    \textbf{Название этапа} & \textbf{Задание} \\
    \hline
    & этапов. 5.~Отчёт необходимо подгрузить в модуле практика как <<письменный отчёт>>. \\
    \hline
    Получение отзыва руководителя практики & Получить отзыв с оценкой у руководителя практики в модуле Практики \\
    \hline
  \end{tabularx}
\end{table}